% results_rq4

\section{RQ4: Bio-Spiritual Inclusion}
\label{sec:rq4}

Islamic law exempts menstruating women from the obligation to pray, and a
prayer tracker that does not represent this exemption misrepresents the
practice it claims to track. RQ4 asks how well the ecosystem serves that need.
Designers do not disagree about how to handle it: almost no one has built it,
the few implementations that exist are not trusted by their own users, and
demand is voiced explicitly and repeatedly across nearly the whole ecosystem.
Inclusion here is an unmet need, not a contested design question.

\subsection{Scarcity Is the Finding}

The \texttt{women\_period} tagger produces 2,194 aspect-tagged sentences,
resolving to 2,182 distinct reviews across the 24 of 26 applications carrying
any (Figure~\ref{fig:rq4-women-aspect}). Its precision is only 20\% (95\% CI
[6, 51]), so we report an adjusted estimate of roughly 439 genuine mentions
(interval [124, 1{,}119]), still spread across nearly the entire ecosystem.
RQ6's promise-and-delivery analysis works from a narrower basis of 1,981
reviews, which is where the sentiment rates below come from; we name all three
units rather than let one stand in for another. Only three applications ---
Athan, Islamic Habit Tracker, and Pillars --- ship any menstrual handling at
all. The rest of this section rests on demand signals and this three-app
scarcity, neither of which depends on the tagger's precision.

\begin{figure}[t]
  \centering
  \includegraphics[width=\linewidth]{fig09_rq4_women_aspect}
  \caption{Reviews mentioning menstruation or prayer exemptions, by application, with negative-sentiment rate. The three applications shipping menstrual handling are named in the text.}
  \label{fig:rq4-women-aspect}
\end{figure}

\subsection{Demand Is Explicit, and Concentrated Where the Feature Is Absent}

Demand-pattern extraction, validated at 91.5\% precision
(\S\ref{sec:method-demand}), is the firm ground here. The 272 demand signals
tagged to \texttt{women\_period} break down into 161 explicit requests, 51
preference citations, 36 churn statements, 21 statements of absence, and 3
tenure citations; 216 come from the 16 applications that both lack the feature
and draw any demand, out of the 23 lacking it. These are signal counts, and one
review can carry several: the same demand resolves to 206 distinct reviews,
the unit RQ6b's unmet-need index uses (\S\ref{sec:rq6}). The register is
direct rather than incidental:

\begin{quote}
\itshape ``a feature for sisters who are on their monthlies would be nice so
they can keep track''
\end{quote}
\hspace*{1em}--- Muslim Pro: Quran Athan Prayer, 4 stars

\subsection{Having the Feature Does Not Improve Sentiment}

If the gap were purely one of coverage, the three applications shipping
menstrual handling should show better sentiment on the topic. They do not:
mentions run 26.3\% negative in those three against 26.0\% in the 21 lacking it
that draw any mention. Shipping the feature is not associated with users being
happier about it, which points at the quality of existing implementations
rather than availability alone:

\begin{quote}
\itshape ``The only complaint I have is: when I turn on menstrual mode it
doesn't turn off the prayers as it promises to do. Please fix it''
\end{quote}
\hspace*{1em}--- Athan: Prayer Times \& Al Quran, 4 stars

\noindent The RQ1 coding corroborates this independently: \emph{menstrual mode
unreliable} appears three times in the single-app pass, and \emph{menstrual
exemption absent} once in the cross-app pass, from a different application
(\S\ref{sec:rq1-coding}). Across two independently drawn samples, this is rare
first-hand evidence that the failure is one of execution as well as coverage.

\subsection{The Between-App Comparison Is Not Estimable}

With three applications on the treated side, no causal or comparative claim
about star rating is possible, and the data show why. The review-weighted
difference in mean star rating is \(+0.085\) (95\% CI \([+0.074, +0.096]\),
computed from the two review pools rather than from M3, whose
\texttt{women\_feature\_score\_diff} term returns neither a standard error nor a
\(p\)-value): apps with the feature average 4.476 stars against 4.391 without
(46,741 reviews across 3 treated apps, 277,946 across 23 untreated).
Averaging first by app and then across apps reverses the sign to \(-0.236\)
(4.135 against 4.371). Both are arithmetically correct and answer different
questions: Athan alone contributes 44,954 of the treated reviews, so the
review-weighted figure is close to one application's rating, while the
unweighted one is pulled down by Islamic Habit Tracker, the corpus's
lowest-rated application. We report \(+0.085\) and label it not estimable
rather than null. A difference reversing direction under a change of
aggregation, on three treated applications, is the clearest demonstration that
\(N = 3\) cannot answer this question.

\subsection{The Topic Model Cannot Add Resolution}

A sub-model fit on women's-health and privacy discourse does not separate
cleanly: after hand-labelling, five of its ten topics carry an identical label
describing one application's data being sold to the US military and a sixth is
a near-duplicate, so six of ten topics sit on a single theme holding 58\% of
the sub-model's non-outlier documents (4,788 of 8,285). One application's
data-handling controversy dominates the subset to the point that the sub-model
cannot further resolve menstrual-handling discourse. We report recurring themes
rather than treat the partition as ten findings, and return to the episode in
\S\ref{sec:limitations}.

